\section{Related work}
\label{sec:related}
\draftnote{the works in this section were located by web search on 7 October 2026 and are described from their abstracts; none has been read in full. A systematic search of IEEE Xplore and Scopus, and the full text of~\cite{markicevic2023ear}, are still required before the statements marked below can stand.}

\paragraph{Efficient super-resolution} Since SRCNN~\cite{dong2016srcnn}, SR networks have grown from a few layers to deep residual and transformer models~\cite{lim2017edsr,ledig2017srgan,liang2021swinir}. A parallel line of work reduces cost under a fixed quality target, with the NTIRE Efficient SR challenges as its reference benchmark~\cite{ren2024ntire,ren2025ntire,ren2026ntire}; RLFN~\cite{kong2022rlfn}, SAFMN~\cite{sun2023safmn}, SMFANet~\cite{zheng2024smfanet} and SPAN~\cite{wan2024span} are representative, and SPAN serves as the challenge baseline. Real-time SR on mobile hardware has its own benchmark series~\cite{ignatov2021mobile}. All of these models are trained and ranked on DIV2K-style data~\cite{agustsson2017div2k} with bicubic downsampling of large natural images. We take their published weights as they are and ask how they behave on a domain, an input size and a degradation the protocol does not cover.

\paragraph{Real-world degradation} That models trained on bicubic downsampling degrade on real images is well established~\cite{lugmayr2019unsupervised}, and has motivated real paired data~\cite{cai2019realsr}, kernel and noise estimation from real images~\cite{bellkligler2019kernelgan,ji2020realsr}, and broad synthetic pipelines that randomise blur, noise, resizing and JPEG compression~\cite{zhang2021bsrgan,wang2021realesrgan}. The latter are usually combined with adversarial training and target perceptual quality on natural images. The quality of the training images themselves also matters: Ohtani et al.~\cite{ohtani2024rethinking} report that training data with few compression artefacts yield better SR models. Closest to our use of real images of a biometric trait is the work of Bulat et al.~\cite{bulat2018learn}, who learn the degradation of real low-resolution faces with a generative network and then train SR on its outputs. Our approach is deliberately simpler, reading the compression level from the files, and our ablation shows that at inputs of 24 to 48 pixels nothing beyond the range of compression levels needs to be estimated.

\paragraph{Super-resolution of compressed images} Joint compression-artefact reduction and SR has been studied for natural images~\cite{dong2015arcnn,jiang2021fbcnn,chen2018cisrdcnn}, has been the subject of a challenge~\cite{yang2022aim}, and has been addressed with efficient networks for streaming content~\cite{ma2024senn}. That SR of compressed inputs requires compression-aware training is therefore not new, and we do not claim it. What we have not found in this literature is a measurement of where published efficient SR models stop being better than interpolation as compression increases, at input sizes of a few dozen pixels, with a check on real images of the target domain \todo{confirm after a systematic search}.

\paragraph{Low resolution and compression in ear recognition} Image resolution is a recognised source of performance variability in unconstrained ear recognition: it is one of the covariates analysed in the Unconstrained Ear Recognition Challenges~\cite{emersic2017uerc,emersic2019uerc}, and Rathgeb et al.~\cite{rathgeb2016compression} study the effect of image compression on ear recognition systems. These studies measure the effect on recognition; they do not ask what an upscaling step does to the image.

\paragraph{Super-resolution for ears and other biometric traits} SR has a long history in face, iris and other biometric modalities~\cite{nguyen2018srbiometrics}. For faces, the benefit of SR for recognising real low-resolution images has been examined directly~\cite{li2019lowres,menezes2021analysis}. For ears, the only study we found is that of Markičević et al.~\cite{markicevic2023ear}, who train EDSR and SwinIR at $\times2$ and $\times4$ on the UERC data and report rank-1 recognition with a ResNet on the super-resolved images, with original and bicubically interpolated images as baselines \todo{state their degradation model and whether compression is considered, after reading the full text}. We differ in the question asked (which published efficient models work on small compressed ear images, and what must change for them to work), in the evaluation (clean ground truth at three sizes, three datasets, paired statistics by subject, and identification on real small images with subjects disjoint from all training), and in the constraint of real-time inference.
